\documentclass{amsart}
\usepackage{amsmath,amssymb,amsthm}
\usepackage{hyperref,cleveref}
\usepackage{tikz,tikz-cd}
\usepackage[all]{xy}

% Theorem-like environments, numbered within sections and sharing a counter.
\newtheorem{theorem}{Theorem}[section]
\newtheorem{lemma}[theorem]{Lemma}
\newtheorem{proposition}[theorem]{Proposition}
\newtheorem{corollary}[theorem]{Corollary}
\theoremstyle{definition}
\newtheorem{definition}[theorem]{Definition}
\newtheorem{example}[theorem]{Example}
\theoremstyle{remark}
\newtheorem*{remark}{Remark}

\numberwithin{equation}{section}

\newcommand{\R}{\mathbb{R}}
\newcommand{\N}{\mathbb{N}}
\newcommand{\norm}[1]{\left\lVert #1 \right\rVert}
\DeclareMathOperator{\supp}{supp}

\title{A Short Tour of \TeX down}
\author{Ada Lovelace \and Bernhard Riemann}
\date{November 2026}

\begin{document}
\maketitle

\begin{abstract}
TeXglass renders \LaTeX{} documents the way a Markdown previewer renders
Markdown. Inline math like $e^{i\pi} + 1 = 0$, displayed equations, and
theorem-like environments are numbered, and every \verb|\ref| resolves to the
right number, even for things defined later (see \cref{thm:main}).
\end{abstract}

\tableofcontents

\section{Preliminaries}\label{sec:prelim}

Throughout, $(X, \mu)$ is a measure space and $1 \le p < \infty$.%
\footnote{The case $p = \infty$ needs the essential supremum instead.}

\begin{definition}[$L^p$ norm]\label{def:norm}
For a measurable $f \colon X \to \R$ we set
\begin{equation}\label{eq:norm}
  \norm{f}_p = \left( \int_X |f|^p \, d\mu \right)^{1/p}.
\end{equation}
\end{definition}

\begin{example}
On $X = \N$ with counting measure, \eqref{eq:norm} becomes
$\norm{x}_p = \bigl(\sum_n |x_n|^p\bigr)^{1/p}$.
\end{example}

\begin{lemma}[Young's inequality]\label{lem:young}
If $a, b \ge 0$ and $\frac1p + \frac1q = 1$ with $p, q > 1$, then
\[
  ab \le \frac{a^p}{p} + \frac{b^q}{q}.
\]
\end{lemma}

\begin{proof}
The logarithm is concave, so
\begin{align}
  \log(ab) &= \tfrac1p \log a^p + \tfrac1q \log b^q \nonumber \\
           &\le \log\Bigl( \frac{a^p}{p} + \frac{b^q}{q} \Bigr). \label{eq:concave}
\end{align}
Exponentiating \eqref{eq:concave} gives the claim.
\end{proof}

\section{Main results}\label{sec:main}

\begin{theorem}[H\"older]\label{thm:main}
Let $p, q$ be as in \cref{lem:young}. Then for all measurable $f, g$,
\begin{equation}\label{eq:holder}
  \int_X |fg| \, d\mu \le \norm{f}_p \norm{g}_q .
\end{equation}
\end{theorem}

\begin{proof}
By homogeneity we may assume $\norm{f}_p = \norm{g}_q = 1$. Apply
Lemma~\ref{lem:young} pointwise with $a = |f(x)|$, $b = |g(x)|$ and integrate.
\end{proof}

\begin{corollary}\label{cor:cs}
Taking $p = q = 2$ in Theorem~\ref{thm:main} gives the Cauchy--Schwarz
inequality $\bigl|\int fg\bigr| \le \norm{f}_2 \norm{g}_2$.
\end{corollary}

\begin{remark}
Unnumbered environments (declared with \verb|\newtheorem*|) work too.
\end{remark}

\begin{proposition}\label{prop:list}
The following hold for $f \in L^p$:
\begin{enumerate}[label=(\roman*)]
  \item\label{it:supp} $\supp f$ is $\sigma$-finite;
  \item\label{it:tri} $\norm{f + g}_p \le \norm{f}_p + \norm{g}_p$ (Minkowski).
\end{enumerate}
\end{proposition}

Item~\ref{it:tri} of \cref{prop:list} follows from \eqref{eq:holder};
see \cite{rudin} for details. \Cref{sec:prelim} collects the definitions,
and \cref{tab:summary} summarises the numbering.

\begin{table}[h]
\centering
\begin{tabular}{|l|c|}
\hline
Object & Number \\
\hline
Definition~\ref{def:norm} & \ref{def:norm} \\
Equation \eqref{eq:holder} & \ref{eq:holder} \\
\hline
\end{tabular}
\caption{Where things ended up.}\label{tab:summary}
\end{table}

\section{Diagrams}\label{sec:diagrams}

Commutative diagrams written with \texttt{tikz-cd} or \texttt{xymatrix} are
drawn in the preview, and can be numbered and referenced like equations. The
square
\begin{equation}\label{eq:square}
\begin{tikzcd}
P \arrow[r, "f"] \arrow[d, "g"', hook] \arrow[dr, phantom, "\lrcorner", very near start] & B \arrow[d, "h", two heads] \\
C \arrow[r, "k"'] & D
\end{tikzcd}
\end{equation}
in \eqref{eq:square} is a pullback.

\begin{lemma}[Snake lemma]\label{lem:snake}
Given a map of short exact sequences
\[
\xymatrix{
0 \ar[r] & A \ar[r]^{f} \ar[d]_{\alpha} & B \ar[r]^{g} \ar[d]^{\beta} & C \ar[r] \ar[d]^{\gamma} & 0 \\
0 \ar[r] & A' \ar[r]_{f'} & B' \ar[r]_{g'} & C' \ar[r] & 0
}
\]
there is an exact sequence $\ker\alpha \to \ker\beta \to \ker\gamma \to \operatorname{coker}\alpha \to \operatorname{coker}\beta \to \operatorname{coker}\gamma$.
\end{lemma}

An adjunction and a natural transformation:
\[
\begin{tikzcd}
\mathcal{C} \arrow[r, shift left=1.2, "F"] & \mathcal{D} \arrow[l, shift left=1.2, "G"]
\end{tikzcd}
\qquad
\begin{tikzcd}[column sep=large]
\mathcal{C} \arrow[r, bend left=45, "F", ""{name=U, below}] \arrow[r, bend right=45, "G"', ""{name=D}] & \mathcal{D}
\arrow[Rightarrow, from=U, to=D, "\eta"]
\end{tikzcd}
\]

Simple TikZ pictures are drawn too:
\begin{center}
\begin{tikzpicture}[scale=1.1]
  \draw[->] (-0.3,0) -- (3.3,0) node[right] {$x$};
  \draw[->] (0,-0.3) -- (0,2.3) node[above] {$y$};
  \draw[thick, blue, domain=0:3.1, samples=40] plot (\x, {0.2*\x*\x});
  \draw[dashed] (2,0) node[below] {$2$} -- (2,0.8) -- (0,0.8) node[left] {$0.8$};
  \fill (2,0.8) circle (1.5pt);
\end{tikzpicture}
\end{center}

\begin{thebibliography}{9}
\bibitem{rudin} W.~Rudin, \emph{Real and Complex Analysis}, McGraw--Hill, 1987.
\end{thebibliography}

\end{document}
